% !PRINTED 299
% ?? uncertain: no folio/running head visible on any of the three bands of this page (section-opening page); number inferred from the printed "300" on the immediately following page.
\weberpage{299}

\begin{center}
\gesperrt{Ninth Section.}

\gesperrt{Estimation of the Roots.}

\rule{0.3\linewidth}{0.4pt}

\S.\,99.

\gesperrt{The Budan-Fourier Theorem.}
\end{center}

The functions of the \gesperrt{Sturm} chains do indeed solve completely and without exception the problem of determining the number of roots of an equation between given limits; but their actual computation is usually difficult and often impracticable. A number of theorems are known for the \gesperrt{estimation} of the number of roots between given limits, which furnish much simpler rules, but which, to be sure, also do not answer the question completely, giving only a maximal number beyond which the number of roots cannot go. These rules are often very useful and sufficient for applications, and they must therefore not be omitted here.

We consider first a procedure which originates with \gesperrt{Budan}, but which was then used and extended by \gesperrt{Fourier}\,\textsuperscript{1}).

Let us consider, in place of a \gesperrt{Sturm} chain, the sequence of derivatives of a real function $f(x)$ of degree $n$

\weq{1}{f(x),\ f'(x),\ f''(x)\ .\ .\ .\ f^{(n)}(x),}

in which, then, $f^{(n)}(x)$ is a constant, which we assume to be different from zero and positive. It is, apart from a positive numerical factor, the coefficient of $x^n$ in $f(x)$.

\noindent\rule{0.25\linewidth}{0.4pt}

\footnotesize
\textsuperscript{1}) \gesperrt{Budan}, Nouvelle methode pour la r\'esolution des \'equations num\'eriques, (presented to the Paris Academy in 1803). \gesperrt{Fourier}, Analyse des \'equations d\'etermin\'ees. Paris 1831. On the history of this question compare \gesperrt{Lagrange}, Trait\'e de la r\'esolution des \'equations, Note VIII (Works, Vol. 8).
\normalsize
